% Generated draft: requires mathematical and source review.
\documentclass[11pt]{article}
\usepackage[margin=25mm]{geometry}
\usepackage{fontspec}
\setmainfont{Latin Modern Roman}
\usepackage{amsmath,amssymb,mathtools,unicode-math}
\setmathfont{Latin Modern Math}
\usepackage{enumitem,needspace,xurl,hyperref}
\setlength{\parindent}{0pt}
\setlength{\parskip}{5pt}
\setlength{\emergencystretch}{3em}
\newcommand{\N}{\mathbb{N}}
\newcommand{\Z}{\mathbb{Z}}
\newcommand{\Q}{\mathbb{Q}}
\newcommand{\R}{\mathbb{R}}
\newcommand{\C}{\mathbb{C}}
\newcommand{\F}{\mathbb{F}}
\newcommand{\A}{\mathbb{A}}
\newcommand{\PP}{\mathbb{P}}
\newcommand{\E}{\mathbb{E}}
\newcommand{\eps}{\varepsilon}
\newcommand{\dd}{\,\mathrm{d}}
\begin{document}
\section*{UnsolvedMath ID: 30001270}
\textbf{Source problem code:} catalogue entry
\subsection{Definitions and mathematical statement}
String topology concerns algebraic operations associated with spaces of loops and paths in a target space. An open--closed topological conformal field theory involves open and closed sectors, with operations associated with suitable moduli spaces of surfaces and compatibility under sewing. These descriptions explain the terminology only; the supplied record does not specify the mathematical models needed to turn them into definitions for this problem.

The supplied question is:
\begin{quote}
Does string topology carry the structure of an open--closed topological conformal field theory?
\end{quote}

The formulation remains unresolved at the level of hypotheses and intended structure. In particular, the record does not specify the target space or manifold, its required hypotheses, coefficient ring, chain-level or homology-level interpretation, open-sector boundary conditions, relevant moduli spaces, or sewing and unit requirements. It also does not state whether the assertion is intended for every target in some class or for a particular target. Consequently, the evidence supports the displayed question, but does not determine a fully quantified existence statement or a unique definition of the requested structure. These choices require source clarification and human mathematical review.
\subsection{Short English statement}
The question asks whether string topology admits compatible open and closed operations forming a topological conformal field theory. The supplied record leaves the target, hypotheses, and precise theory unspecified, so a fully precise formulation cannot be recovered from this evidence alone.
\subsection{Sources}
\begin{itemize}
\item Supplied local UnsolvedMath registry record: permanent ID \(30001270\), exact source problem code \texttt{catalogue entry}, in the Oberwolfach Reports Open Problems collection. The record lists proposed year \(2009\) and supplies no individual proposer.
\item Supplied report citation: \emph{Strings, Fields and Topology} (2009), Oberwolfach Reports, pp.\ 1547--1596, DOI \texttt{10.4171/owr/2009/28}. Source URL: \url{https://doi.org/10.4171/owr/2009/28}. The same URL is supplied under the literature-source label ``original OWR report.'' No problem-specific page, speaker attribution, or passage defining the intended structure is supplied.
\item Literature assessment dated \texttt{2026-08-21}: this is supplied dated curation, not an independently verified assessment. It labels the problem open because, according to that assessment, the extraction contains only the word ``conjecture'' and no mathematical statement, so no resolution could be verified. The background describes this as dated literature triage and requests independent verification of time-sensitive status and cited claims.
\item The supplied statement metadata labels the statement \texttt{corrected\_verified} and says that the mathematical question was preserved while open--closed typography was normalized. This supplied verification label and the supplied open-status label are not independent mathematical or source verification by this draft.
\end{itemize}
\end{document}
